\begin{abstract}
Molecular-machine-learning scores depend on how molecules are represented, grouped, and assigned to validation partitions. We present a paired chemometric audit of one benchmark-construction intervention: selecting scaffold-disjoint test sets with lower train--test target-mean mismatch while holding test-set size and candidate-search budget fixed. Six public molecular-property datasets were rebuilt with explicit molecular-identity and duplicate policies. For each of 20 pre-specified partition seeds, a target-blind scaffold candidate pool was generated; a size-matched baseline was selected without target information, and a lower-target-mean-gap counterpart was selected only among candidates with exactly the same test size. Partition manifests, hashes, search budgets, and inferential units were frozen before model fitting.

No classification dataset--model cell met the pre-specified corrected criterion for a target-balanced advantage. All six primary regression cells favored target-mean-aware selection under the convention that all acyclic molecules share one scaffold identity, but these effects attenuated strongly when acyclic molecules were treated as singleton scaffold groups. A mean-only control was added to distinguish learned-model improvement from the direct coupling between target-mean mismatch and RMSE difficulty. A dominant-fragment sensitivity left all nine affected classification inferences inconclusive while reversing seven of nine point-estimate directions.

The results show that benchmark effects are conditional on partition-selection rules, scaffold semantics, and molecular-record representation. Molecular benchmark performance should therefore be interpreted as protocol-conditioned evidence rather than a model-intrinsic quantity. The complete audit uses exact-size paired perturbations, frozen search budgets, molecule-level manifests, partition-level inference, and predeclared sensitivity analyses.
\end{abstract}

\noindent\textbf{Keywords:} chemometrics; molecular property prediction; benchmark validation; scaffold splitting; reproducibility
